\documentclass[aspectratio=169,10pt]{beamer}
\usetheme{Madrid}
\usecolortheme{default}

\usepackage{hyperref}
\usepackage{tikz}
\usepackage{amsmath,amssymb}

% --- Custom Colors ---
\definecolor{unibblue}{RGB}{0,51,102}
\definecolor{unibgold}{RGB}{212,175,55}
\definecolor{unibgreen}{RGB}{0,128,0}

\setbeamercolor{palette primary}{bg=unibblue,fg=white}
\setbeamercolor{palette secondary}{bg=unibblue!85,fg=white}
\setbeamercolor{palette tertiary}{bg=unibblue!70,fg=white}
\setbeamercolor{structure}{fg=unibblue}
\setbeamercolor{block title}{bg=unibblue!85,fg=white}
\setbeamercolor{block body}{bg=unibblue!10,fg=black}
\setbeamercolor{alertblock title}{bg=unibgold!90,fg=black}
\setbeamercolor{alertblock body}{bg=unibgold!15,fg=black}

\title{Persamaan Diferensial}
\subtitle{Minggu 15: Sintesis PDP dalam Persamaan Maxwell \& Panduan Team-Based Project}
\author{Novalio Daratha \& Muhammad Arfan}
\institute{Program Studi Teknik Elektro\\Universitas Bengkulu}
\date{2026}

\begin{document}

\begin{frame}
  \titlepage
\end{frame}

\begin{frame}{Capaian Pembelajaran (CPMK 4 \& C6)}
  \begin{itemize}
    \item \textbf{CPMK 4.1:} Mengintegrasikan seluruh konsep PDB dan PDP ke dalam 4 Persamaan Maxwell fundamental.
    \item \textbf{CPMK 4.2:} Menurunkan Persamaan Gelombang Elektromagnetik 3D ($\nabla^2 \mathbf{E} = \mu\varepsilon \frac{\partial^2 \mathbf{E}}{\partial t^2}$) dari hukum Faraday dan Ampere-Maxwell.
    \item \textbf{CPMK 4.3:} Menghubungkan kecepatan gelombang elektromagnetik dengan konstanta ruang hampa ($c = 1/\sqrt{\mu_0\varepsilon_0}$).
    \item \textbf{CPMK 4.4 (C6):} Menyajikan hasil sintesis pemodelan matematika dalam bentuk laporan dan presentasi *Team-Based Project*.
  \end{itemize}
\end{frame}

\begin{frame}{Daftar Isi}
  \tableofcontents
\end{frame}

\section{Empat Persamaan Diferensial Maxwell}
\begin{frame}{Empat Persamaan Maxwell dalam Bentuk Diferensial}
  \begin{block}{1. Hukum Gauss untuk Medan Listrik}
    $$ \nabla \cdot \mathbf{D} = \rho_v \quad \iff \quad \text{Muatan listrik adalah sumber divergensi medan listrik} $$
  \end{block}
  \pause
  \begin{block}{2. Hukum Gauss untuk Medan Magnetik}
    $$ \nabla \cdot \mathbf{B} = 0 \quad \iff \quad \text{Tidak ada monopol magnetik (garis medan magnet selalu kontinu/tertutup)} $$
  \end{block}
  \pause
  \begin{block}{3. Hukum Induksi Faraday}
    $$ \nabla \times \mathbf{E} = -\frac{\partial \mathbf{B}}{\partial t} \quad \iff \quad \text{Perubahan medan magnetik menghasilkan medan listrik pusaran} $$
  \end{block}
  \pause
  \begin{block}{4. Hukum Ampere-Maxwell}
    $$ \nabla \times \mathbf{H} = \mathbf{J} + \frac{\partial \mathbf{D}}{\partial t} \quad \iff \quad \text{Arus konduksi dan arus pergeseran menghasilkan medan magnetik} $$
  \end{block}
\end{frame}

\section{Penurunan Persamaan Gelombang EM 3D}
\begin{frame}{Dari Persamaan Maxwell ke Persamaan Gelombang 3D}
  Pada ruang hampa bebas muatan ($\rho_v = 0, \mathbf{J} = 0$):
  $$ \nabla \times \mathbf{E} = -\mu_0 \frac{\partial \mathbf{H}}{\partial t}, \quad \nabla \times \mathbf{H} = \varepsilon_0 \frac{\partial \mathbf{E}}{\partial t} $$
  \pause
  Terapkan operator $\nabla \times (\nabla \times \mathbf{E})$ dan identitas vektor $\nabla \times (\nabla \times \mathbf{A}) = \nabla(\nabla \cdot \mathbf{A}) - \nabla^2 \mathbf{A}$:
  $$ \nabla(\underbrace{\nabla \cdot \mathbf{E}}_{=0}) - \nabla^2 \mathbf{E} = -\mu_0 \frac{\partial}{\partial t}(\nabla \times \mathbf{H}) = -\mu_0 \frac{\partial}{\partial t}\left(\varepsilon_0 \frac{\partial \mathbf{E}}{\partial t}\right) $$
  \pause
  \begin{alertblock}{Persamaan Gelombang Elektromagnetik 3D}
    $$ {\color{unibblue} \nabla^2 \mathbf{E} = \mu_0 \varepsilon_0 \frac{\partial^2 \mathbf{E}}{\partial t^2}} \quad \iff \quad \nabla^2 \mathbf{E} = \frac{1}{c^2} \frac{\partial^2 \mathbf{E}}{\partial t^2} $$
    dengan $c = \frac{1}{\sqrt{\mu_0 \varepsilon_0}} \approx 3 \times 10^8\text{ m/s}$ (Kecepatan Cahaya).
  \end{alertblock}
\end{frame}

\section{Panduan Team-Based Project}
\begin{frame}{Pedoman Team-Based Project (Bobot 20\%)}
  \begin{block}{Ketentuan Proyek Kelompok (3--4 Mahasiswa)}
    Pilihlah salah satu topik pemodelan riil teknik elektro:
    \begin{enumerate}
      \item Analisis Stabilitas Transien Sistem Tenaga Listrik (*Swing Equation - Nonlinear ODE*).
      \item Perambatan Sinyal Frekuensi Tinggi pada Jaringan 5G (*Telegrapher PDE*).
      \item Distribusi Arus Efek Kulit pada Konduktor Tegangan Ekstra Tinggi (*Bessel ODE*).
      \item Pemodelan Distribusi Medan Potensial Gardu Induk 2D (*Laplace PDE*).
    \end{enumerate}
  \end{block}
  \pause
  \begin{alertblock}{Format Luaran (Deliverables)}
    \begin{itemize}
      \item Laporan Ilmiah LaTeX (Maksimal 10 Halaman).
      \item Kode Simulasi Komputasi Julia / Pluto Notebook.
      \item Presentasi Slide Beamer (10 menit presentasi + 5 menit tanya jawab).
    \end{itemize}
  \end{alertblock}
\end{frame}

\begin{frame}{Rubrik Penilaian Team-Based Project}
  \begin{table}[]
    \centering
    \small
    \begin{tabular}{|l|c|l|}
      \hline
      \textbf{Kriteria Penilaian} & \textbf{Bobot} & \textbf{Deskripsi} \\ \hline
      Ketepatan Pemodelan Fisik & 25\% & Penurunan PDB/PDP dari hukum dasar elektro. \\ \hline
      Keakuratan Solusi Analitik & 25\% & Penurunan matematis langkah-demi-langkah. \\ \hline
      Kualitas Simulasi Komputasi & 25\% & Validasi numerik menggunakan Julia / visualisasi grafik. \\ \hline
      Kualitas Laporan \& Presentasi & 25\% & Sistematika penulisan LaTeX dan kejelasan presentasi. \\ \hline
    \end{tabular}
  \end{table}
\end{frame}

\begin{frame}{Penutup Mata Kuliah}
  \begin{center}
    {\Large \textbf{Selamat Mempersiapkan Ujian Akhir Semester (UAS)!}}\\
    \vspace{0.8cm}
    \textit{"Persamaan diferensial adalah bahasa alami untuk mendeskripsikan alam semesta dan seluruh inovasi teknologi rekayasa elektro."}\\
    \vspace{1cm}
    \textbf{Tim Dosen Pengampu:}\\
    Novalio Daratha, Ph.D. \& Muhammad Arfan, M.T.\\
    Program Studi Teknik Elektro, Universitas Bengkulu
  \end{center}
\end{frame}

\end{document}
